\section{Notation}\label{sec:notation}

\sumpara The three tables list every symbol once, with its meaning and units. \sumparaend
\fin \Cref{tab:notation-data,tab:notation-schedule,tab:notation-method} list every symbol in the report with its meaning, its units, and where it is defined. \fine
\fin No symbol carries two meanings. \fine

\begin{table}[!htbp]
\centering
\caption{Indices and data.}
\label{tab:notation-data}
\small
\begin{tabularx}{\textwidth}{lLll}
\toprule
Symbol & Meaning & Units & Defined in \\
\midrule
$n$ & Number of jobs & jobs & \cref{sec:timing} \\
$m$ & Number of stages & stages & \cref{sec:timing} \\
$j$ & Job index, $j = 1, \ldots, n$ & & \cref{sec:timing} \\
$h$ & Job index in sums and maxima that already involve $j$ & & \cref{sec:mip} \\
$k$ & Stage index in line order, $k = 1, \ldots, m$ & & \cref{sec:timing} \\
$i$ & Position index within a sequence $\pi$ & & \cref{sec:timing} \\
$q$ & Position index in the mixed-integer program & & \cref{sec:mip} \\
$r_j$ & Release time of job $j$ & minutes & \cref{sec:timing} \\
$d_j$ & Due date of job $j$ & minutes & \cref{sec:timing} \\
$w_j$ & Weight (priority score) of job $j$ & integer score & \cref{sec:timing} \\
$p_{jk}$ & Processing time of job $j$ at stage $k$ & minutes & \cref{sec:timing} \\
\bottomrule
\end{tabularx}
\end{table}

\begin{table}[!htbp]
\centering
\caption{Schedules, measures, and variables of the mixed-integer program.}
\label{tab:notation-schedule}
\small
\begin{tabularx}{\textwidth}{lLll}
\toprule
Symbol & Meaning & Units & Defined in \\
\midrule
$\pi$, $\pi(i)$ & A sequence of all jobs, and the job in its position $i$ & & \cref{sec:timing} \\
$S_{jk}$, $C_{jk}$ & Start and completion of job $j$ at stage $k$ & minutes & \eqref{eq:recursion} \\
$T_j$ & Tardiness of job $j$, $\max\{0, C_{jm} - d_j\}$ & minutes & \cref{sec:timing} \\
$Z(\pi)$ & Total weighted tardiness of $\pi$ (primary measure) & weighted minutes & \eqref{eq:measures} \\
$W(\pi)$ & Total waiting between stages of $\pi$ (secondary measure) & minutes & \eqref{eq:measures} \\
$F(\pi)$ & The pair $(Z(\pi), W(\pi))$, compared lexicographically & & \cref{sec:methods} \\
$x_{jq}$ & 1 if job $j$ takes position $q$, 0 otherwise & binary & \cref{sec:mip} \\
$B_{qk}$, $E_{qk}$ & Start and end of the job in position $q$ at stage $k$ & minutes & \cref{sec:mip} \\
$U_j$ & Upper bound on the tardiness of job $j$ in \eqref{eq:mip-tardy} & minutes & \cref{sec:mip} \\
\bottomrule
\end{tabularx}
\end{table}

\begin{table}[!htbp]
\centering
\caption{Method and experiment settings.}
\label{tab:notation-method}
\small
\begin{tabularx}{\textwidth}{lLll}
\toprule
Symbol & Meaning & Value or units & Defined in \\
\midrule
$\pi^{\mathrm{EDD}}$ & Sequence of the earliest due date rule & & \cref{alg:ig} \\
$\pi^{*}$, $\pi'$, $\pi''$ & Best sequence found, a rebuilt sequence, and a local search candidate & & \cref{alg:ig} \\
$\delta$ & Number of jobs removed in each iteration & 4 & \cref{sec:methods} \\
$\lambda$ & Temperature factor & 0.5 & \cref{sec:methods} \\
$\Theta$ & Temperature, $\lambda \sum_{j,k} p_{jk} / (10\, n\, m)$, compared with increases in $Z$ & minutes & \cref{sec:methods} \\
$\Delta$ & Increase in $Z$ of a candidate over the current sequence & weighted minutes & \cref{sec:methods} \\
$L$, $\ell$ & Iteration limit and iteration counter & $L = 5000$ & \cref{alg:ig} \\
$u$ & Uniform random draw on $[0, 1]$ in the acceptance test & & \cref{alg:ig} \\
$\omega$ & Scenario index & 30 per test & \cref{sec:stress-design} \\
$\varepsilon_{jk\omega}$ & Lognormal factor on $p_{jk}$ in scenario $\omega$, mean 1 & & \cref{sec:stress-design} \\
$\nu$ & Coefficient of variation of each factor & 0.1 to 0.8 & \cref{sec:stress-design} \\
$\sigma$ & Standard deviation of $\ln \varepsilon_{jk\omega}$, $\sigma^2 = \ln(1 + \nu^2)$ & & \cref{sec:stress-design} \\
\bottomrule
\end{tabularx}
\end{table}
